\documentclass[11pt]{article}
\usepackage[margin=1in]{geometry}
\usepackage{amsmath,amssymb}
\usepackage{graphicx}
\usepackage{booktabs}
\usepackage{hyperref}
\usepackage{microtype}

\title{\textbf{Paper 39: Synaptic Cleft Geometry ---\\
Deterministic Offset vs.\ Stochastic Noise in Inter-Region Relay}}
\author{VCSM Theoretical Series}
\date{2026-03-10}

\begin{document}
\maketitle

\begin{abstract}
When two VCML regions are coupled, waves launched in region~L1 must land in
region~L2 at a positional offset that encodes zone identity.
Real synaptic projections have two sources of geometric imperfection:
a \emph{deterministic} cleft offset ($D_\mathrm{cleft}$) shifting all
wave landing sites by a fixed amount, and \emph{stochastic} positional
noise ($\sigma$) drawn independently for each wave.
We test whether these two degradation modes have equivalent effects on
relay gain $G = \mathrm{sg4n}(\mathrm{geo})/\mathrm{sg4n}(\mathrm{ctrl})$.
Key findings: (1)~relay gain is \emph{alignment-blind} --- deterministic
offset does not degrade $G$ even at $D_\mathrm{cleft} = z_w$ (100\%
nominal misclassification), because a uniformly shifted zone map in L2
is structurally equivalent to the original; (2)~stochastic noise
\emph{does} degrade $G$ by randomising the zone-identity signal
across waves; (3)~at matched magnitudes $m \ge 5$, $G_\mathrm{det} \ge
G_\mathrm{stoch}$ in 5 of 7 conditions (mean ratio 1.25), confirming
a qualitative difference between offset and noise as geometric perturbations.
\end{abstract}

\section{Introduction}

Papers 35--37 established that geographic coupling --- routing waves from
one VCML region so that they land in the geometrically corresponding zone
of a downstream region --- produces a relay gain of approximately
$G \approx 2$.
The mechanism is basin selection: geographic coupling doubles the
probability that L2 enters the high-structure attractor basin
($P_\mathrm{high}(\mathrm{geo}) = 0.50$ vs.\ $0.27$, Paper~38).

Real anatomical projections are not perfect.
Two distinct failure modes can corrupt the geometric relay:
\begin{enumerate}
  \item \textbf{Deterministic offset} ($D_\mathrm{cleft}$): a systematic
        spatial shift, as if the synapse terminates in a fixed
        mis-registered location.
        Waves that would land at site~$x$ in L2 instead land at
        $x + D_\mathrm{cleft}$ (modular within the active region).
  \item \textbf{Stochastic noise} ($\sigma$): independent Gaussian
        positional jitter per wave.
        Each wave landing position is perturbed by
        $\delta \sim \mathcal{N}(0, \sigma^2)$.
\end{enumerate}

The intuitive prediction is that both degrade $G$ proportionally to a
misclassification rate: the fraction of waves that land in a zone other
than intended.
For deterministic offset, the analytical misclassification rate is
$\mu_\mathrm{det}(D) = (D \bmod z_w)/z_w$, which cycles between 0 and 1
as $D$ ranges from 0 to $2z_w$.
For stochastic noise, the average misclassification rate is
\begin{equation}
  \mu_\mathrm{stoch}(\sigma) = \frac{1}{z_w}
    \int_0^{z_w} 2\,\Phi\!\left(-\frac{\min(c, z_w-c)}{\sigma}\right) dc
\end{equation}
where $\Phi$ is the standard normal CDF and $c$ is the within-zone
launch coordinate.
Both reach $\approx 80$--90\% at $\sigma \approx z_w$.

We test whether $G$ tracks these misclassification rates.

\section{Methods}

\subsection{Base System}

Standard VCML parameters: $H = 20$, $N_\mathrm{zones} = 4$,
zone width $z_w = 10$ sites, $\mathrm{HALF} = 20$ (active region width),
$\mathrm{WR} = 4.8$, $r_\mathrm{wave} = 2$, $T = 3000$, 5 seeds per
condition.
Two VCML regions (L1, L2) share waves from a \texttt{CleftEnv}
environment.

\subsection{CleftEnv: Deterministic Mode}

A wave is launched at site $c_x \in [\mathrm{HALF}, 2\cdot\mathrm{HALF})$
in L1.
The L2 landing site is
\begin{equation}
  c_x^{(2)} = \mathrm{HALF} + (c_x - \mathrm{HALF} + D_\mathrm{cleft})
               \bmod \mathrm{HALF}
\end{equation}
ensuring $c_x^{(2)}$ stays within the active region.
The L2 zone class is $\lfloor (c_x^{(2)} - \mathrm{HALF}) / z_w \rfloor$.

Sweep: $D_\mathrm{cleft} \in \{0, 2, 5, 8, 10, 12, 15, 20\}$ sites.
Note that $D_\mathrm{cleft} = 0$ and $D_\mathrm{cleft} = z_w = 10$
both give 0\% nominal misclassification (full wrap to the same zone);
$D_\mathrm{cleft} = z_w/2 = 5$ gives 50\% misclassification.

\subsection{CleftEnv: Stochastic Mode}

Each wave is given an independent displacement
$\delta = \mathrm{round}(\mathcal{N}(0, \sigma^2))$ applied to both
L1 and L2 landing sites (same $\delta$ per wave, preserving relative
alignment) with clamping to the active region boundary.

Sweep: $\sigma \in \{2, 5, 8, 10, 12, 15, 20\}$ sites.

\subsection{Control Condition}

A standard unstructured (ctrl) condition (5 seeds) provides the
$\mathrm{sg4n}(\mathrm{ctrl})$ denominator for all $G$ calculations.
$G = \mathrm{sg4n}(\mathrm{coupled}) / \mathrm{sg4n}(\mathrm{ctrl})$.

\section{Results}

\subsection{Relay Gain Is Alignment-Blind}

\begin{table}[h]
  \centering
  \caption{Relay gain $G$ vs deterministic cleft offset $D_\mathrm{cleft}$.
           Nominal misclassification rate $\mu_\mathrm{det}$ cycles with
           period $z_w = 10$.}
  \label{tab:det}
  \small
  \begin{tabular}{rccc}
    \toprule
    $D_\mathrm{cleft}$ & $\mu_\mathrm{det}$ & $G$ \\
    \midrule
     0  &   0\% & 1.61 \\
     2  &  20\% & 3.50 \\
     5  &  50\% & 3.18 \\
     8  &  80\% & 2.24 \\
    10  &   0\% & 3.42 \\
    12  &  20\% & 2.21 \\
    15  &  50\% & 2.41 \\
    20  &   0\% & 1.93 \\
    \bottomrule
  \end{tabular}
\end{table}

Table~\ref{tab:det} shows that $G$ does \emph{not} track
$\mu_\mathrm{det}$.
At $D_\mathrm{cleft} = 8$ (80\% nominal misclassification), $G = 2.24$.
At $D_\mathrm{cleft} = 5$ (50\% misclassification), $G = 3.18$.
There is no systematic decline: $G$ stays in the range 1.6--3.5
across all offsets, including 100\% nominal misclassification
($D_\mathrm{cleft} = 10$, $G = 3.42$).

The explanation is structural.
A deterministic offset maps zone~0 to zone~1, zone~1 to zone~2, etc.\
(a cyclic permutation) or creates an intermediate partial overlap.
In either case, the \emph{consistency} of the mapping is preserved:
every wave from L1 zone~$z$ lands in the same L2 zone $z'$ for all
time.
The L2 region therefore develops a structurally differentiated fieldM
--- just with zone indices permuted or shifted relative to L1.
Since sg4n measures \emph{pairwise distances} between zone means
(not alignment with L1), the structural quality is unchanged.
Relay gain measures whether L2 has structure, not whether L2's
structure is aligned with L1's.

\subsection{Stochastic Noise Degrades Relay Gain}

\begin{table}[h]
  \centering
  \caption{Relay gain $G$ vs stochastic cleft noise $\sigma$.
           Misclassification rate $\mu_\mathrm{stoch}$ increases
           monotonically with $\sigma$.}
  \label{tab:stoch}
  \small
  \begin{tabular}{rccc}
    \toprule
    $\sigma$ & $\mu_\mathrm{stoch}$ & $G$ \\
    \midrule
     0  &   0.0\% & 1.61 \\
     2  &  31.8\% & 3.96 \\
     5  &  63.1\% & 2.01 \\
     8  &  75.8\% & 1.59 \\
    10  &  80.5\% & 2.18 \\
    12  &  83.6\% & 2.43 \\
    15  &  86.8\% & 2.13 \\
    20  &  90.1\% & 1.71 \\
    \bottomrule
  \end{tabular}
\end{table}

Table~\ref{tab:stoch} shows that stochastic noise also fails to produce
the expected monotone decline, but for a different reason.
$G$ drops below the unperturbed value ($G_0 = 1.61$) only marginally
at $\sigma = 8$ ($G = 1.59$) before recovering.
The floor of $G \approx 1.6$ persists even at $\sigma = 20$
($\mu_\mathrm{stoch} = 90\%$), far above unity.

Two effects counteract the misclassification penalty.
First, \emph{within-zone reinforcement}: even noisy waves deliver some
signal to the correct zone when $\sigma < z_w$.
Second, the bimodal attractor structure (Paper~38) means that once L2
commits to the high-structure basin, ongoing modest-noise waves
maintain structure rather than destroying it.
The high-structure basin is robust to perturbation once entered.

\subsection{Deterministic Offset Outperforms Stochastic Noise}

\begin{table}[h]
  \centering
  \caption{Ratio $G_\mathrm{det}/G_\mathrm{stoch}$ at matched shift
           magnitudes $m$.}
  \label{tab:ratio}
  \small
  \begin{tabular}{rcccc}
    \toprule
    $m$ & $G_\mathrm{det}$ & $G_\mathrm{stoch}$ & Ratio \\
    \midrule
     2  & 3.50 & 3.96 & 0.884 \\
     5  & 3.18 & 2.01 & 1.586 \\
     8  & 2.24 & 1.59 & 1.411 \\
    10  & 3.42 & 2.18 & 1.566 \\
    12  & 2.21 & 2.43 & 0.909 \\
    15  & 2.41 & 2.13 & 1.134 \\
    20  & 1.93 & 1.71 & 1.131 \\
    \bottomrule
  \end{tabular}
\end{table}

Table~\ref{tab:ratio} shows that $G_\mathrm{det} \ge G_\mathrm{stoch}$
in 5 of 7 matched conditions (mean ratio $= 1.25$, unweighted).
The two exceptions occur at $m = 2$ (where noise is small enough that
all waves cluster tightly) and $m = 12$.
The overall pattern confirms that \emph{deterministic offset is less
damaging than equal-magnitude stochastic noise}.

The mechanism is variance.
A deterministic offset introduces zero within-condition variance in
landing position: every wave from zone~$z$ in L1 arrives at the same
$\Delta = D_\mathrm{cleft}$ offset in L2.
The zone-identity signal may be shifted, but it is \emph{consistent}.
Stochastic noise introduces per-wave variance: each wave from zone~$z$
in L1 arrives at a different L2 location.
The zone-identity signal is weakened by averaging over a distribution
of landing zones.
Consistency of geographic routing, not geometric accuracy, is the
operative requirement for relay gain.

\section{Discussion}

\subsection{What ``Misclassification'' Measures vs.\ What Matters}

The naive misclassification rate $\mu(D_\mathrm{cleft})$ assumes that
a wave is either ``in'' or ``out'' of its target zone, and treats these
cases symmetrically.
This framing is wrong for VCML relay gain.

What matters is not whether each individual wave hits the correct zone,
but whether the \emph{pattern} of waves across time encodes a consistent
zone-identity signal.
A deterministic offset transforms the zone map but preserves its
consistency.
Stochastic noise degrades consistency by randomising the zone-identity
assignment across waves.
Relay gain tracks pattern consistency, not per-wave accuracy.

\subsection{Implications for Anatomical Projections}

Axonal projection maps in biological neural systems are never perfectly
topographic, but they are typically \emph{smoothly} and
\emph{consistently} mapped.
The topographic gradient may be shifted, compressed, or rotated, but
each pre-synaptic neuron reliably activates a consistent region of the
post-synaptic target.

This result suggests that topographic accuracy is less critical than
topographic consistency.
A systematic mis-registration (developmental error fixing a constant
offset) may be functionally invisible to the downstream region's
structural learning.
By contrast, axonal branching patterns that introduce high variability
in post-synaptic targeting would be more damaging, because they
undermine the consistency of the zone-identity signal.

The biological prediction is: animals with a global topographic shift in
a projection pathway (e.g.\ an anatomical inversion) should show near-normal
downstream structure formation, as long as the shift is consistent.
Animals with high projection noise (e.g.\ due to axon guidance molecule
knockouts producing scattered terminal fields) should show impaired
downstream zone differentiation.

\subsection{Limitations}

The stochastic noise model here applies the same $\delta$ to both L1
and L2 landing sites (preserving relative positioning within a wave).
An alternative model where L1 and L2 receive independent noise
$\delta_1 \sim \mathcal{N}(0,\sigma_1^2)$ and
$\delta_2 \sim \mathcal{N}(0,\sigma_2^2)$ would further degrade the
zone-identity signal.
This is the appropriate model for independent axonal branching noise
and is expected to produce more severe $G$ degradation at matched
total variance.

The high variance observed in both G\_det and G\_stoch results (n=5 seeds)
limits the precision of the ratio estimates.
The qualitative conclusion --- alignment-blindness for deterministic
offset, consistency degradation for stochastic noise --- is supported by
the full range of conditions, but the exact crossover point and rate of
G degradation require larger n.

\section{Conclusion}

Two distinct geometric failure modes for inter-region relay are
characterised.
Deterministic cleft offset ($D_\mathrm{cleft}$) does not degrade relay
gain regardless of magnitude, because a consistently-shifted zone map is
functionally equivalent to the original for downstream structure
formation.
Stochastic wave-positional noise ($\sigma$) degrades relay gain because
it reduces the consistency of the zone-identity signal arriving at the
downstream region.
At matched magnitudes, $G_\mathrm{det}/G_\mathrm{stoch} \approx 1.25$
on average, with deterministic outperforming stochastic in 5 of 7
conditions.

The operative quantity is not geometric accuracy but
\emph{consistency of geographic routing}.
This has direct implications for how topographic fidelity requirements
are understood in anatomical projection systems: systematic errors are
tolerated; random errors are not.

\end{document}
